% GENERATED FILE. Edit canonical lessons, then run npm run generate:books.
% canonical-source-hash: fnv1a64:c45753d16333d745
% canonical-lessons: RU-C219-zubnaya-pasta, RU-C219-britva, RU-C219-utyug, RU-C219-pylesos, RU-C219-dukhovka

\chapter{Toothpaste, Razor, Iron, Vacuum cleaner, Oven}
\label{ch:ru-toothpaste-razor-iron-vacuum-cleaner-oven}
\hlchaptermodality{219}

\begin{chapteropening}
\textbf{By the end of this chapter:} \emph{I can say toothpaste, a razor, an iron, a vacuum cleaner and an oven.}

Complete the last lesson of chapter 219: I can say toothpaste, a razor, an iron, a vacuum cleaner and an oven.
\end{chapteropening}

\section[\texorpdfstring{\ru{зубная} \ru{паста} (zubnáya pásta)}{зубная паста (zubnáya pásta)}]{\ru{зубная} \ru{паста} (zubnáya pásta) — toothpaste}

\label{lesson:RU-C219-zubnaya-pasta}



\begin{quote}
Before the new one: say the Russian for main, then the Russian for ordinary.
\end{quote}

\subsection*{You'll want to know: \ru{зубная} \ru{паста}}
\textbf{\ru{зубная} \ru{паста}} (\emph{zubnáya pásta}) — \textquotedblleft{}toothpaste\textquotedblright{}.

\begin{grammarlens}[title={What you've built}]
One more word for this chapter.
\end{grammarlens}

\subsection*{Your turn}
\begin{itemize}
\raggedright
  \item \emph{Say it:} \emph{zubnáya pásta}
  \item \emph{Say it:} \emph{zubnáya pásta}, once more
  \item \emph{Say it:} say \emph{zubnáya pásta}
  \item \emph{Recall:} say the Russian for main, then the Russian for ordinary, then say \emph{zubnáya pásta} again
\end{itemize}

\begin{tcolorbox}[breakable,colback=teal!4,colframe=teal!35!black,title={Before you move on}]
What does \ru{зубная} \ru{паста} mean? (\textquotedblleft{}Toothpaste\textquotedblright{}.) Say it once more.
\end{tcolorbox}
\section[\texorpdfstring{\ru{бритва} (brítva)}{бритва (brítva)}]{\ru{бритва} (brítva) — a razor}

\label{lesson:RU-C219-britva}



\begin{quote}
Before the new one: say the Russian for ordinary, then the Russian for toothpaste.
\end{quote}

\subsection*{You'll want to know: \ru{бритва}}
\textbf{\ru{бритва}} (\emph{brítva}) — \textquotedblleft{}a razor\textquotedblright{}.

\begin{grammarlens}[title={What you've built}]
One more word for this chapter.
\end{grammarlens}

\subsection*{Your turn}
\begin{itemize}
\raggedright
  \item \emph{Say it:} \emph{brítva}
  \item \emph{Say it:} \emph{brítva}, once more
  \item \emph{Say it:} say \emph{brítva}
  \item \emph{Recall:} say the Russian for ordinary, then the Russian for toothpaste, then say \emph{brítva} again
\end{itemize}

\begin{tcolorbox}[breakable,colback=teal!4,colframe=teal!35!black,title={Before you move on}]
What does \ru{бритва} mean? (\textquotedblleft{}A razor\textquotedblright{}.) Say it once more.
\end{tcolorbox}
\section[\texorpdfstring{\ru{утюг} (utyúg)}{утюг (utyúg)}]{\ru{утюг} (utyúg) — an iron}

\label{lesson:RU-C219-utyug}



\begin{quote}
Before the new one: say the Russian for toothpaste, then the Russian for a razor.
\end{quote}

\subsection*{You'll want to know: \ru{утюг}}
\textbf{\ru{утюг}} (\emph{utyúg}) — \textquotedblleft{}an iron\textquotedblright{}.

\begin{grammarlens}[title={What you've built}]
One more word for this chapter.
\end{grammarlens}

\subsection*{Your turn}
\begin{itemize}
\raggedright
  \item \emph{Say it:} \emph{utyúg}
  \item \emph{Say it:} \emph{utyúg}, once more
  \item \emph{Say it:} say \emph{utyúg}
  \item \emph{Recall:} say the Russian for toothpaste, then the Russian for a razor, then say \emph{utyúg} again
\end{itemize}

\begin{tcolorbox}[breakable,colback=teal!4,colframe=teal!35!black,title={Before you move on}]
What does \ru{утюг} mean? (\textquotedblleft{}An iron\textquotedblright{}.) Say it once more.
\end{tcolorbox}
\section[\texorpdfstring{\ru{пылесос} (pylesós)}{пылесос (pylesós)}]{\ru{пылесос} (pylesós) — a vacuum cleaner}

\label{lesson:RU-C219-pylesos}



\begin{quote}
Before the new one: say the Russian for a razor, then the Russian for an iron.
\end{quote}

\subsection*{You'll want to know: \ru{пылесос}}
\textbf{\ru{пылесос}} (\emph{pylesós}) — \textquotedblleft{}a vacuum cleaner\textquotedblright{}.

\begin{grammarlens}[title={What you've built}]
One more word for this chapter.
\end{grammarlens}

\subsection*{Your turn}
\begin{itemize}
\raggedright
  \item \emph{Say it:} \emph{pylesós}
  \item \emph{Say it:} \emph{pylesós}, once more
  \item \emph{Say it:} say \emph{pylesós}
  \item \emph{Recall:} say the Russian for a razor, then the Russian for an iron, then say \emph{pylesós} again
\end{itemize}

\begin{tcolorbox}[breakable,colback=teal!4,colframe=teal!35!black,title={Before you move on}]
What does \ru{пылесос} mean? (\textquotedblleft{}A vacuum cleaner\textquotedblright{}.) Say it once more.
\end{tcolorbox}
\section[\texorpdfstring{\ru{духовка} (dukhóvka)}{духовка (dukhóvka)}]{\ru{духовка} (dukhóvka) — an oven}

\label{lesson:RU-C219-dukhovka}



\begin{quote}
Before the new one: say the Russian for an iron, then the Russian for a vacuum cleaner.
\end{quote}

\subsection*{You'll want to know: \ru{духовка}}
\textbf{\ru{духовка}} (\emph{dukhóvka}) — \textquotedblleft{}an oven\textquotedblright{}.

\begin{grammarlens}[title={What you've built}]
One more word for this chapter.
\end{grammarlens}

\subsection*{Your turn}
\begin{itemize}
\raggedright
  \item \emph{Say it:} \emph{dukhóvka}
  \item \emph{Say it:} \emph{dukhóvka}, once more
  \item \emph{Say it:} say \emph{dukhóvka}
  \item \emph{Recall:} say the Russian for an iron, then the Russian for a vacuum cleaner, then say \emph{dukhóvka} again
\end{itemize}

\begin{tcolorbox}[breakable,colback=teal!4,colframe=teal!35!black,title={Before you move on}]
What does \ru{духовка} mean? (\textquotedblleft{}An oven\textquotedblright{}.) Say it once more.
\end{tcolorbox}
